\documentclass[11pt]{article}
\usepackage[margin=1in]{geometry}
\usepackage{amsmath,amssymb,amsthm,mathtools}
\usepackage[hidelinks]{hyperref}
\newtheorem{theorem}{Theorem}
\newtheorem{lemma}[theorem]{Lemma}
\newtheorem{corollary}[theorem]{Corollary}
\title{A real-rooted companion and backward certificates for gap-three sections}
\author{Internal proof and exact replay; no priority claim}
\date{17 November 2026}
\begin{document}
\maketitle
The \href{GAP_THREE_WRONSKIAN_RECURRENCE.tex}{Wronskian proof} reduces
simplicity of $T_{n,\gamma,3}=\sum_k\binom nk(\gamma)_{3k}t^k$ to
coprimality of three real cyclic sections $W_0,W_1,W_2$.
A nonzero real-rooted section suffices, but its uniform existence is
unproved. We give two precise related structures, retaining the missing
transfer and arithmetic boundaries. General FC and unrestricted SFC
remain open.

\section{A finite-difference representation and its real-rooted companion}
Let $B=B_{n,\gamma,3}$ be the monic degree-$n$ auxiliary, defined by
\[
 B_{0,\gamma}=1,\qquad
 B_{n+1,\gamma}=(z+3n+\gamma+2)B_{n,\gamma}-zB_{n,\gamma}'.
\]
Put $p_n(x)=\prod_{j=0}^{n-1}(x+2+3j)$ and
$\nabla f(x)=f(x)-f(x-1)$.
\begin{lemma}\label{lem:diff}
For every real $\gamma$,
\[
 e^{-z}B(z)=\sum_{k\ge0}p_n(\gamma-k)\frac{(-z)^k}{k!},
 \qquad B(z)=\sum_{j=0}^n\frac{\nabla^jp_n(\gamma)}{j!}z^j.
                                                               \tag{1}
\]
\end{lemma}
\begin{proof}
Writing $\delta=zD_z$, the recurrence gives
$e^{-z}B=\prod_{j=0}^{n-1}(\gamma+2+3j-\delta)e^{-z}$.
On the $k$th Taylor monomial $\delta$ acts by $k$, proving the first
identity as an identity of entire functions (or formal series).
Multiplication by $e^z$ gives the finite backward differences in (1).
Terms above degree $n$ vanish because $p_n$ has degree $n$.
\end{proof}

Write $B(z)=\sum_{j=0}^n b_jz^j$, and define
\[
 J(y)=\sum_{k=0}^n\binom nk p_n(\gamma-k)y^k,
 \qquad K(z)=\sum_{j=0}^n n_{\underline j}b_jz^j.              \tag{2}
\]
\begin{theorem}\label{thm:polar}
For every $n\ge1$ and $\gamma>n-2$, both $J$ and $K$ have $n$
distinct real roots, all strictly less than $-1$.
In particular this holds throughout the needed range
$\gamma=n+\beta-1$, $\beta>0$.
\end{theorem}
\begin{proof}
Let $q$ have degree $n$ and only negative real roots, allowing
multiplicities, and let $c>n$. The polynomial $(c-yD_y)q$ retains a
root of multiplicity $m-1$ at each root of multiplicity $m$ of $q$.
All remaining roots are simple, one strictly to the left of the least
root and one in each interval between consecutive distinct roots.
Indeed, away from roots its zeros solve
\[
 c-\sum_r\frac{m_r y}{y-r}=0.
\]
The left side has strictly negative derivative
$\sum_r m_r r/(y-r)^2<0$. Its endpoint limits give exactly the stated
roots; on the rightmost interval it has no zero since its limit at
positive infinity is $c-n>0$. No newly created root coincides with an
old root, and leading and constant coefficients remain nonzero.

Now
\[
 J(y)=\prod_{j=0}^{n-1}(\gamma+2+3j-yD_y)(1+y)^n.
\]
All $n$ constants are greater than $n$. The preceding rule shows that
the original root $-1$ loses one multiplicity at each step, while every
new root is simple and lies to its left. After $n$ steps all roots are
simple and strictly below $-1$.

The binomial transform of (1) gives
\[
 J(y)=(1+y)^nK\left(-\frac{y}{1+y}\right).                   \tag{3}
\]
Neither endpoint degenerates: $J(-1)=\nabla^np_n(\gamma)=n!$
and $K(-1)=p_n(\gamma-n)>0$. The real map
$y\mapsto-y/(1+y)$ sends $(-\infty,-1)$ bijectively to itself,
with nonzero derivative. This proves the assertion for $K$ as well.
\end{proof}

There is a related published finite-difference theorem: Iskander--Jain,
\emph{Hyperbolicity of Appell polynomials of functions in the
$\delta$-Laguerre--P\'olya class}, Theorem 1.1, in the
\href{sources/iskander-jain-2020-delta-laguerre-polya.pdf}{cached arXiv version}
(\href{https://arxiv.org/abs/2002.09906}{primary source}).
Its finite zero-free interval applies to $p_n$ on $[\gamma-n,\gamma]$.
The elementary polar proof above is independent of that input and also
proves simplicity and the location below $-1$.
The source's Theorem 1.2 has a DIFFERENT infinite-ray hypothesis.
It cannot be applied to the backward Poisson series in (1): the sampled
polynomial $p_n(\gamma-s)$ has positive roots $\gamma+2+3j$, so its
forward ray from zero is not zero-free. No hyperbolicity theorem for $B$
or its Wronskian sections follows from that theorem here.

\section{The companion property alone does not transfer to the sections}
The real-rootedness of $K$ does not, for an arbitrary inverse transform,
guarantee even one real-rooted individual Wronskian section. An exact
control uses $n=8$ and
\[
 K_*(z)=\prod_{j=0}^7\left(1+\frac{90+j}{108}z\right),\qquad
 B_*(z)=\sum_{j=0}^8\frac{[z^j]K_*}{8_{\underline j}}z^j.
\]
Rescale $B_*$ to be monic and form $W$ by the earlier gap-three
Wronskian formula. $K_*$ has eight distinct roots below $-1$, but the
three sections of $W$ have respectively degrees $5,5,4$ and only
$3,3,0$ real roots, by exact Sturm counts.
This control is NOT the arithmetic-progression Gamma family. It refutes
only a generic inverse-companion inference to an individual section,
not a theorem exploiting the full factors $\gamma+2+3j$, nor possible
real-rooted linear combinations, nor SFC.

\section{A backward certificate needs only nonzero pivots}
\begin{theorem}\label{thm:pivots}
Let $N\ge3$ and let $q_N,q_{N-1},q_{N-2}$ be real polynomials with
constant coefficient one and degree at most $\lfloor j/3\rfloor$ at
index $j$. Starting at $j=N$, put
\[
 \lambda_j=[u](q_j-q_{j-1}),\qquad
 q_{j-3}=\frac{q_j-q_{j-1}}{\lambda_ju},                     \tag{4}
\]
and continue downward through $j=3$.
If every pivot $\lambda_j$ is nonzero, then the input triple has gcd
one over $\mathbb C$.
No sign condition on the pivots, root-location theorem, or claim of
real-rootedness of the input polynomials is needed.
\end{theorem}
\begin{proof}
The numerator has zero constant coefficient and its linear coefficient
is $\lambda_j$. Thus each new $q_{j-3}$ is polynomial with constant
coefficient one. Its degree bound follows by division by $u$.
Consequently $q_0=q_1=q_2=1$ at termination.
If a nonzero $u_0$ were a common root of a current triple, (4) would
make it a root of the next polynomial. Descending all the way to
$q_0=1$ is impossible. Zero is not a common root since constants are one.
\end{proof}
For the actual moment family at $n\ge2$, take $N=2n$ and
$q_{2n-r}=W_r/W_r(0)$, $r=0,1,2$.
The constants are nonzero in the physical range $\gamma=n+\beta-1$,
$\beta>0$, by the proved coefficient positivity. Completion of this
backward certificate therefore proves simplicity, independently of the
stronger real-rooted-section strategy. Uniform nonvanishing of these
pivots has NOT been proved.

For positive pivots, the same recurrence has an elementary weighted
interval interpretation: $q_N$ sums $u^{|S|}\prod_{j\in S}\lambda_j$
over $S\subset\{3,\ldots,N\}$ whose distinct elements differ by at least
three. This is the partition polynomial of nonoverlapping length-three
intervals on a path. The identity follows by including or excluding the
last interval; its interpretation is not needed for the gcd proof.
A zero pivot means only that this particular backward certificate is
inconclusive. For example the full-degree positive input
$q_6=1+u+u^2$, $q_5=1+u$, $q_4=1+2u$ has gcd one but first pivot zero.
It is not permissible to label such an index a failure of SFC.

\section{Native-prime residuals also require a further step}
At the ordinary shape $\gamma=n$, use
$T_{n,n,3}(t)=\sum_k\binom nk(n)_{3k}t^k$.
The assertion that all its Newton residuals at three are squarefree is
false: at $n=22$ the edge $(19,27)$ to $(22,33)$ has residual
$1+Y^3=(1+Y)^3$ in $\mathbb F_3[Y]$.
Parity alone does not repair the native-prime test uniformly.
At $n=63$ the edges $(54,80)$ to $(60,89)$ and $(60,89)$ to $(63,95)$
have residuals $1+2Y^3$ and $1+Y^3$.
The first has reduced slope denominator two and permits degree two
for a hypothetical square divisor. Thus the residual-and-parity test
still permits a quadratic. These are exact limitations of sufficient
certificates, not repeated-root or SFC counterexamples.

The checker pins the source, proves the polynomial identities symbolically,
replays positive-pivot certificates, and checks the stated counterprofiles.
The larger-index Sturm controls at $n=24,48$ are diagnostics only.
No uniform pivot theorem, uniform real-rooted-section existence,
all-index gap-three SFC, general FC or unrestricted SFC is claimed.
All novelty, independent-replay, formal-verification and external-review
flags remain false.
\end{document}
